The intervention defines the reach of a reliance claim. Strong responses to zero vectors alongside weak responses to reassignment or a training-derived centroid suggest sensitivity to representation support. Singleton groups introduce a complementary boundary: their original content-position relations never enter an entity-preserving reassignment. Coverage and control choice therefore determine which dependence has been tested, including when the measured effect is small.

The construct controls show why whole-content dependence needs a second level of interpretation. A source label is meaningful language, yet dependence on that label differs from dependence on the directional proposition. Selective edits distinguish these stipulated cues, while paraphrase effects expose sensitivity to how the same generated truth is represented. This distinction makes cue-specific evidence useful without elevating it to unrestricted semantic identification.

Capability evidence addresses the value of the evaluated response. Improvement over persistence establishes a different comparison from improvement over a numerical/timing model, while removing the complete text stream measures another quantity. These comparators separate overall prediction, incremental content benefit and complete-stream utility (Table~\ref{tab:capability_overview}). A perturbation effect remains defined when capability is weak or unassessed. Interpreting a small effect as evidence that useful context is absent, however, requires a capable predictor and an effective intervention.
